% ---------------------------------------------------------------
% seccoes/representacao.tex — Representação do conhecimento
% ---------------------------------------------------------------

\chapter{Representação do conhecimento}\label{sec:4}

Neste capítulo é apresentada a forma como o conhecimento adquirido foi representado no sistema pericial, através da definição dos factos, eventos e regras utilizados pelo motor de inferência.

\section{Dos sensores aos factos}\label{sec:4.1}

Cada sensor simulado tem uma cadência e um alcance próprios, que determinam a frequência e a distância a que a informação chega às regras (Tabela~\ref{tab:5}).

\begin{tabela}{L{\dimexpr0.2195\linewidth-2\tabcolsep\relax}L{\dimexpr0.1298\linewidth-2\tabcolsep\relax}L{\dimexpr0.2496\linewidth-2\tabcolsep\relax}L{\dimexpr0.3992\linewidth-2\tabcolsep\relax}}
\caption{Sensores simulados}\label{tab:5}\\
\linhacab \cab{\textbf{Sensor}} & \cab{\textbf{Cadência}} & \cab{\textbf{Alcance}} & \cab{\textbf{Informação para as regras}} \\
\endfirsthead
\linhacab \cab{\textbf{Sensor}} & \cab{\textbf{Cadência}} & \cab{\textbf{Alcance}} & \cab{\textbf{Informação para as regras}} \\
\endhead
Radar frontal & 100 ms & 130 m & Distância, velocidade relativa e TTC do veículo da frente; veículos à frente nas vias vizinhas \\
\rowcolor{corFundo}Câmara & 200 ms & 120 m × visibilidade (84 m com chuva) & Peões, sinais, semáforos, buracos, obstáculos, barreiras e marcas da via \\
Radares de ângulo morto & 200 ms & 30 m (perigos até 60 m) & Ocupação de cada via lateral \\
\rowcolor{corFundo}Sensor de velocidade & 100 ms & — & Velocidade do próprio veículo \\
Sensor de ambiente & 400 ms & 200 m & Chuva, visibilidade e aderência \\
\end{tabela}

A informação dos sensores é passada para as regras de três formas distintas:

\begin{itemize}
  \item medições diretas, nomeadamente, distância, velocidade;
  \item dados derivados, como o TTC e a distância necessária para parar;
  \item \textbf{factos semânticos}, como o risco crítico e a via ocupada, que resultam da comparação de um dado derivado com um critério.
\end{itemize}

\section{Tipos de conhecimento}\label{sec:4.2}

Para decidir, o sistema precisa também de saber como esses factos evoluem no tempo, que regras são aplicadas, que critérios devem ser utilizados quando a norma não é suficiente e como resolver conflitos entre decisões. Neste sentido, o conhecimento foi organizado em seis tipos, apresentados na Tabela~\ref{tab:6}.

\begin{tabela}{L{\dimexpr0.2182\linewidth-2\tabcolsep\relax}L{\dimexpr0.7798\linewidth-2\tabcolsep\relax}}
\caption{Tipos de conhecimento representados}\label{tab:6}\\
\linhacab \cab{\textbf{Tipo}} & \cab{\textbf{Exemplos no sistema}} \\
\endfirsthead
\linhacab \cab{\textbf{Tipo}} & \cab{\textbf{Exemplos no sistema}} \\
\endhead
Factual & Velocidade atual, veículo à frente, estado de um semáforo, ocupação de uma via \\
\rowcolor{corFundo}Temporal & Três leituras de TTC baixo em 2 s; paragem completa de 1,2 s no sinal STOP; TTC baixo registado em simultâneo com chuva \\
Normativo & Parar no vermelho e no STOP, ceder passagem a peões, respeitar o limite, circular pela direita, ultrapassar pela esquerda \\
\rowcolor{corFundo}Heurístico & Curvas de travagem para peões e obstáculos; desvio para o lado provado seguro; desistir de uma ultrapassagem impossível \\
Difuso & Velocidade adequada para situações de visibilidade reduzida ou piso escorregadio \\
\rowcolor{corFundo}Metaconhecimento & Ordem das fases (perceção, decisão, recurso); categorias de prioridade; critérios de desempate \\
\end{tabela}

O conhecimento factual e o temporal descrevem a situação, ou seja, o que está a acontecer e se se mantém ao longo de várias leituras. O conhecimento normativo é proveniente do Código da Estrada, e o heurístico vem sobretudo do perito e da literatura. Por sua vez, o conhecimento difuso possibilita que os conceitos, como a visibilidade reduzida, não sejam definidos por uma fronteira rígida e o metaconhecimento, não decide nenhuma manobra, uma vez que é utilizado para definir a ordem pela qual as regras são avaliadas e qual das opções propostas prevalece quando várias são ativadas ao mesmo tempo. Importa ainda distinguir a base de conhecimento, que contém as regras, da memória de trabalho, que mantém os factos e estados que o sistema conhece num dado momento.

\section{Estrutura das regras e regras representativas}\label{sec:4.3}

Todas as regras de decisão seguem a mesma estrutura: \textbf{SE} condições sobre factos e eventos \textbf{ENTÃO} propor uma ação, com velocidade recomendada, nível de risco, prioridade, razão e evidência. Uma regra nunca atua sobre o veículo; apenas propõe uma ação.

\begin{caixaregra}
SE    condição 1\\
E     condição 2\\
E     condição 3\\
ENTÃO ação proposta\\
\hspace*{3.2em}NÍVEL DE RISCO = nível\\
\hspace*{3.2em}PRIORIDADE DA DECISÃO = categoria
\end{caixaregra}

A base de conhecimento, atualmente, tem cerca de 123 regras e vários ficheiros DRL, organizados por situação como, por exemplo, peões, semáforos, veículo à frente, mudança de via e meteorologia. A Tabela~\ref{tab:7} apresenta algumas das regras representativas do projeto

\begin{tabela}{L{\dimexpr0.0695\linewidth-2\tabcolsep\relax}L{\dimexpr0.3171\linewidth-2\tabcolsep\relax}L{\dimexpr0.2314\linewidth-2\tabcolsep\relax}L{\dimexpr0.1267\linewidth-2\tabcolsep\relax}L{\dimexpr0.1038\linewidth-2\tabcolsep\relax}L{\dimexpr0.1496\linewidth-2\tabcolsep\relax}}
\caption{Regras representativas da base de conhecimento}\label{tab:7}\\
\linhacab \cab{\textbf{\#}} & \cab{\textbf{SE}} & \cab{\textbf{ENTÃO}} & \cab{\textbf{Risco}} & \cab{\textbf{Prior.}} & \cab{\textbf{Tipo}} \\
\endfirsthead
\linhacab \cab{\textbf{\#}} & \cab{\textbf{SE}} & \cab{\textbf{ENTÃO}} & \cab{\textbf{Risco}} & \cab{\textbf{Prior.}} & \cab{\textbf{Tipo}} \\
\endhead
R1 & Velocidade acima do limite em vigor & Reduzir até ao limite & Moderado & 20 & Normativo \\
\rowcolor{corFundo}R2 & Chuva, visibilidade ou aderência reduzidas & Reduzir para a velocidade calculada por lógica difusa & Moderado ou alto & 40 & Difuso \\
R3 & Semáforo vermelho e veículo junto à linha & Parar e manter parado & Alto & 60 & Normativo \\
\rowcolor{corFundo}R4 & Sinal STOP, em qualquer via, a menos de 4 m da linha & Parar 1,2 s e depois avançar & Alto & 60 & Normativo e temporal \\
R5 & Peão no caminho ou a atravessar, entre 6 e 32 m & Parar & Alto & 60 & Normativo \\
\rowcolor{corFundo}R6 & Peão no caminho ou a atravessar, a 6 m ou menos & Travagem de emergência & Crítico & 70 & Heurístico \\
R7 & Pelo menos 3 leituras de TTC < 2,5 s nos últimos 2 s & Travagem de emergência & Crítico & 70 & Temporal \\
\rowcolor{corFundo}R8 & Chuva e pelo menos 2 leituras de TTC < 3,0 s em 2 s, cada uma coincidente com uma leitura de chuva & Travar até 15 km/h & Alto & 50 & Temporal \\
R9 & Obstáculo na via a 50 m ou menos e mudança para a esquerda provada segura & Mudar para a esquerda & Baixo & 50 & Heurístico \\
\rowcolor{corFundo}R10 & Pedido de mudança de via com a via de destino ocupada & Manter a via & Moderado & 30 & Normativo \\
\end{tabela}

\section{Níveis de risco e prioridades}\label{sec:4.4}

O perito sugeriu quatro níveis (baixo, médio, alto e crítico). Durante a implementação verificou-se que um único eixo não chegava para ordenar propostas de natureza diferente: uma regra de conforto e uma regra que evita uma colisão podem ambas ter risco «moderado». Separaram-se, por isso, dois conceitos: os quatro níveis passaram a descrever o risco da situação (Figura~\ref{fig:3}), e foi criada uma prioridade com sete categorias para decidir entre propostas (Tabela~\ref{tab:8}). Desta forma, o risco acompanha a decisão e a explicação, mas não intervém na seleção.

\begin{figure}[H]
\centering
\includegraphics[width=0.81\linewidth,height=0.80\textheight,keepaspectratio]{figuras/fig03_niveis_de_risco}
\caption{Níveis de risco}\label{fig:3}
\end{figure}

\begin{tabela}{L{\dimexpr0.2688\linewidth-2\tabcolsep\relax}L{\dimexpr0.0924\linewidth-2\tabcolsep\relax}L{\dimexpr0.2886\linewidth-2\tabcolsep\relax}L{\dimexpr0.3482\linewidth-2\tabcolsep\relax}}
\caption{Categorias de prioridade e o seu papel na arbitragem}\label{tab:8}\\
\linhacab \cab{\textbf{Categoria}} & \cab{\textbf{Valor}} & \cab{\textbf{Papel na arbitragem da velocidade}} & \cab{\textbf{Exemplos}} \\
\endfirsthead
\linhacab \cab{\textbf{Categoria}} & \cab{\textbf{Valor}} & \cab{\textbf{Papel na arbitragem da velocidade}} & \cab{\textbf{Exemplos}} \\
\endhead
Segurança e emergência & 70 & Sobrepõe-se a tudo & Peão a 6 m ou menos; TTC crítico \\
\rowcolor{corFundo}Paragem legal obrigatória & 60 & Sobrepõe-se a tudo abaixo de 70 & Semáforo vermelho; STOP; peão a atravessar \\
Prevenção da colisão & 50 & Teto de velocidade & Veículo à frente; obstáculo; TTC baixo com chuva \\
\rowcolor{corFundo}Perigo rodoviário & 40 & Teto de velocidade & Buraco; meteorologia (lógica difusa) \\
Segurança na manobra & 30 & Base, se nenhum teto for inferior & Arrancar após paragem; retomar no verde; mudanças de via \\
\rowcolor{corFundo}Cumprimento da velocidade & 20 & Teto de velocidade & Limite em vigor e limite à frente; não ultrapassar pela direita \\
Progressão e conforto & 10 & Base, se nada mais se aplicar & Manter a velocidade; pedido de aceleração \\
\end{tabela}
